\section*{Capítulo \ref{ch:g11:subduction-zones} --- Las zonas de subducción}

\begin{solution}{exo:g11:subduction-zones:1}
La \omterm{def:g11:subduction-zones:subduction}{subducción} es el hundimiento de una placa de \omterm{def:g11:ocean-ridges:oceanic-lithosphere}{litosfera oceánica} bajo otra
placa, de vuelta al manto; una \omterm{def:g11:subduction-zones:subduction}{zona de subducción} es un límite de placas
convergente donde se produce, señalado por una fosa. Se hunde la \omterm{def:g11:ocean-ridges:oceanic-lithosphere}{litosfera
oceánica}: cuando es antigua, es fría, gruesa y más densa que la \omterm{def:g11:lithosphere-asthenosphere:lithosphere}{astenosfera}
que tiene debajo. La corteza continental es demasiado ligera para hundirse.
\end{solution}

\begin{solution}{exo:g11:subduction-zones:2}
(1) Por ejemplo, la región de Tonga--Fiyi, el Pacífico occidental de Japón a
las Marianas, Indonesia y América del Sur. (2) Siempre del lado de la placa
cabalgante, nunca del lado oceánico. (3) Someros: todos los terremotos de la
dorsal mesoatlántica del mapa están a menos de \qty{70}{km} de profundidad.
\end{solution}

\begin{solution}{exo:g11:subduction-zones:3}
$7031 \times 0.21 \approx 1480$ terremotos a más de \qty{70}{km};
$7031 \times 0.069 \approx 485$ a más de \qty{300}{km}.
\end{solution}

\begin{solution}{exo:g11:subduction-zones:4}
El \omterm{def:g11:subduction-zones:eclogite}{esquisto azul} es una roca metamórfica formada a partir del basalto de la
\omterm{def:g11:ocean-ridges:oceanic-lithosphere}{corteza oceánica} enterrado a alta presión mientras sigue relativamente frío,
coloreada de azul por la glaucofana; la \omterm{def:g11:subduction-zones:eclogite}{eclogita} es la roca metamórfica densa
en la que se convierten el basalto y el \omterm{def:g11:ocean-ridges:gabbro}{gabro} a las profundidades del manto,
formada por granate rojo y piroxeno verde. Las dos proceden de la \omterm{def:g11:ocean-ridges:oceanic-lithosphere}{corteza
oceánica} y conservan su composición química; lo que cambia son los minerales,
formados en estado sólido con la nueva presión y la nueva temperatura, y la
densidad.
\end{solution}

\begin{solution}{exo:g11:subduction-zones:5}
Fondo oceánico: \omterm{def:g11:lithosphere-asthenosphere:lithosphere}{litosfera} antigua que se mueve hacia la fosa. Fosa: la placa se
dobla hacia abajo y empieza a hundirse. Antearco: se acumulan los sedimentos
arrancados de la placa que se hunde; debajo, las dos placas se deslizan una
contra otra en el límite de placas, donde se producen grandes terremotos. Arco
volcánico: una línea de volcanes sobre la \omterm{def:g11:subduction-zones:slab}{losa}. Trasarco: detrás de algunos
arcos, la placa cabalgante se estira.
\end{solution}

\begin{solution}{exo:g11:subduction-zones:6}
$\tan\delta = (508 - 135)/(800 - 300) = 373/500 \approx 0.746$, luego
$\delta \approx \ang{37}$.
\end{solution}

\begin{solution}{exo:g11:subduction-zones:7}
(1) $\tan\ang{45} = 1$: a \qty{300}{km} de profundidad. (2) A \qty{670}{km} de
la fosa. (3) $670\sqrt{2} \approx \qty{948}{km}$ de \omterm{def:g11:subduction-zones:slab}{losa}.
\end{solution}

\begin{solution}{exo:g11:subduction-zones:8}
(1) $(74.6 + 70.0 + 65.5 + 63.1 + 52.6)/5 \approx \qty{65.2}{mm}$ al año;
$(47.6 + 31.0 + 29.2 + 22.8 + 19.8 + 15.2 + 14.9 + 10.6)/8 \approx
\qty{23.9}{mm}$ al año; cociente $65.2/23.9 \approx 2.7$. (2) Las placas se
mueven más deprisa cuando de ellas cuelga una \omterm{def:g11:subduction-zones:slab}{losa}: el \omterm{def:g11:subduction-zones:slab-pull}{tirón de la losa} que se
hunde es la fuerza más intensa que mueve las placas, como concluyeron Forsyth
y Uyeda.
\end{solution}

\begin{solution}{exo:g11:subduction-zones:9}
(1) \qty{83}{mm} son \qty{0.083}{m}: $50/0.083 \approx 600$ años y
$60/0.083 \approx 720$ años de movimiento de la placa. (2) De 869 a 2011 van
1142 años: del mismo orden. Si el límite de placas permanece bloqueado entre
grandes terremotos, el desplazamiento de 2011 liberó varios siglos de
movimiento acumulado.
\end{solution}

\begin{solution}{exo:g11:subduction-zones:10}
(1) Japón: $6400 - 3200\,e^{-133/62.8} \approx \qty{6015}{m}$; Tonga:
$6400 - 3200\,e^{-99/62.8} \approx \qty{5740}{m}$; Andes:
$2500 + 350\sqrt{49} = \qty{4950}{m}$. (2) El fondo que llega a Japón: el más
antiguo, y por tanto la \omterm{def:g11:lithosphere-asthenosphere:lithosphere}{litosfera} más fría y más densa, y el fondo más hondo.
\end{solution}

\begin{solution}{exo:g11:subduction-zones:11}
(1) $3380 \times 1.06 \approx \qty{3583}{kg/m^3}$. (2)
$(3583 - 2900)/2900 \approx 0.24$: un aumento de alrededor del \qty{24}{\percent}.
(3) Antes, la corteza es más ligera que el manto en
$3380 - 2900 = \qty{480}{kg/m^3}$; después, es más densa en unos
\qty{203}{kg/m^3}. Ya no sostiene la \omterm{def:g11:subduction-zones:slab}{losa}: tira de ella hacia abajo.
\end{solution}

\begin{solution}{exo:g11:subduction-zones:12}
(1) $t = 6 \times 10^6 \times 3.16 \times 10^7 \approx 1.9 \times 10^{14}$~s;
$\sqrt{10^{-6} \times 1.9 \times 10^{14}} \approx \qty{1.4e4}{m}$, unos
\qty{14}{km}. (2) Mucho menos que \qty{100}{km}: en 6 millones de años solo se
calienta la piel exterior de la \omterm{def:g11:subduction-zones:slab}{losa}; su interior sigue frío. (3) La roca fría
es rígida y puede romperse, mientras que el manto caliente que la rodea fluye;
los terremotos profundos se producen dentro de la \omterm{def:g11:subduction-zones:slab}{losa} fría.
\end{solution}

\begin{solution}{exo:g11:subduction-zones:13}
(1) $\sqrt{1017^2 + 670^2} \approx \qty{1218}{km}$. (2) \qty{64}{mm} al año
son \qty{64}{km} por millón de años: $1218/64 \approx 19$ millones de años.
(3) Unas tres veces más que bajo Tonga: la \omterm{def:g11:subduction-zones:slab}{losa} bajo los Andes se hunde más
despacio y con un buzamiento más suave.
\end{solution}

\begin{solution}{exo:g11:subduction-zones:14}
(1) Una rebanada de \qty{1}{m} de ancho a lo largo de la fosa tiene un volumen
de $10^5 \times 10^6 \times 1 = \qty{1e11}{m^3}$; su exceso de masa es
$36 \times 10^{11} = \qty{3.6e12}{kg}$. (2) Su exceso de peso es
$3.6 \times 10^{12} \times 9.8 \approx \qty{3.5e13}{N}$ por cada metro de
fosa. (3) La \omterm{def:g11:subduction-zones:slab}{losa} está unida al resto de la placa, que es rígida: la fuerza se
transmite a toda la placa, que es arrastrada hacia la fosa, tanto más cuanto
que las fosas se extienden a lo largo de miles de kilómetros.
\end{solution}

\begin{solution}{exo:g11:subduction-zones:15}
(1) $(23.6 - 16.5) \times 111 \approx \qty{790}{km}$. (2) La velocidad
aumenta $237 - 104 = \qty{133}{mm}$ al año en \qty{790}{km}, unos
\qty{17}{mm} al año por cada cien kilómetros. (3) La convergencia es la
velocidad de la placa del Pacífico menos la del arco. Como la velocidad de la
placa del Pacífico es más o menos la misma en los dos puntos, el propio arco
tiene que moverse hacia la placa del Pacífico, cada vez más deprisa hacia el
norte: el trasarco que tiene detrás se está estirando, más en el norte que en
el sur.
\end{solution}

\begin{solution}{pb:g11:subduction-zones:1}
\textbf{1.} A unos \qty{671}{km}.

\textbf{2.} Del lado oceánico la placa todavía no se ha hundido: allí la
\omterm{def:g11:lithosphere-asthenosphere:lithosphere}{litosfera} está cerca de la superficie. Los terremotos profundos se producen
dentro de la \omterm{def:g11:subduction-zones:slab}{losa}, que solo está bajo la placa cabalgante.

\textbf{3.} $z(300) = -24 + 273 = \qty{249}{km}$ y
$z(700) = -24 + 637 = \qty{613}{km}$; $\tan\delta = 364/400 = 0.91$, luego
$\delta \approx \ang{42}$.

\textbf{4.} $x = (100 + 24)/0.91 \approx \qty{136}{km}$ y
$x = (400 + 24)/0.91 \approx \qty{466}{km}$.

\textbf{5.} Una recta con este buzamiento que empezara en la fosa pasaría por
$z = 0$ en $x = 0$; la recta queda por encima. La \omterm{def:g11:subduction-zones:slab}{losa} no es recta: cerca de la
fosa buza más suavemente, y se empina en profundidad.

\textbf{6.} \qty{16.8}{cm} al año, \qty{168}{km} por millón de años.

\textbf{7.} $168/69 \approx 2.4$ veces más rápida.

\textbf{8.} La convergencia a ambos lados de la fosa es la velocidad de la
placa del Pacífico menos la velocidad del arco del otro lado. Si el arco se
mueve hacia la placa del Pacífico, como ocurre mientras se estira el trasarco
que tiene detrás, las dos velocidades se suman, y la convergencia es mayor que
la velocidad de la placa del Pacífico sola.

\textbf{9.} $6400 - 3200\,e^{-99/62.8} \approx \qty{5740}{m}$.

\textbf{10.} Al cabo de unos 60 millones de años, la \omterm{def:g11:ocean-ridges:oceanic-lithosphere}{litosfera oceánica} es, en
su conjunto, más densa que la \omterm{def:g11:lithosphere-asthenosphere:lithosphere}{astenosfera} que tiene debajo; a los 99 millones
de años es todavía más densa. En cuanto su borde desciende, su peso arrastra
al resto.

\textbf{11.} $(670 + 24)/0.91 \approx \qty{763}{km}$.

\textbf{12.} $\sqrt{763^2 + 670^2} \approx \qty{1015}{km}$.

\textbf{13.} $1015/168 \approx 6.0$ millones de años.

\textbf{14.} $1015/69 \approx 14.7$ millones de años.

\textbf{15.} $3380 \times 1.06 \approx \qty{3583}{kg/m^3}$: la corteza, antes
más ligera que el manto, pasa a ser más densa que él, y toda la \omterm{def:g11:subduction-zones:slab}{losa} tira con
más fuerza.

\textbf{16.} $t \approx 6.0 \times 10^6 \times 3.16 \times 10^7 \approx
1.9 \times 10^{14}$~s, y $\sqrt{\kappa t} \approx \qty{14}{km}$.

\textbf{17.} \qty{14}{km} es poco al lado de \qty{100}{km}: solo se ha
calentado la parte exterior de la \omterm{def:g11:subduction-zones:slab}{losa}; su interior sigue frío cuando llega a
\qty{670}{km}.

\textbf{18.} La \omterm{def:g11:subduction-zones:slab}{losa} fría sigue siendo lo bastante rígida para romperse, hasta
su extremo, mientras el manto caliente que la rodea fluye: los terremotos bajo
Tonga se producen dentro de la \omterm{def:g11:subduction-zones:slab}{losa}, hasta unos \qty{671}{km}.

\textbf{19.} Las ondas sísmicas muestran un cambio brusco en el manto a unos
\qty{670}{km}: allí sus velocidades dan un salto.

\textbf{20.} Unos \textbf{6 millones de años} desde la fosa hasta
\qty{670}{km}.
\end{solution}
